\documentstyle[% 


righttag% 


,11pt% 


,amssymb% 


,russian%               remove in Eng. 


,izvru%                 Set izven% in Eng. 


]{amsart} 


 


% 


 


\newcommand{\deq}{\overset{\mathrm{d}}{=}{}} 


\newcommand{\tts}[1]{} 


 


%\hoffset=-15mm%                  For Eng. 


  


\setcounter{page}{83} 


\god{1999} 


\nomer{4~(443)} %                        Remove in Eng. 


%\nomer{Vol.\,43, No.\,4}%               Set in Eng. 


%\str{\pageref{begin}--\pageref{end}}%  Set in Eng. 


\udk{519.216} 


 


\author{\it A.Н.\,Чупрунов} 


% NB:   ^^ remove in Eng. 


\title{Принцип инвариантности для независимых наблюдений 


банаховозначного случайного процесса} 


 


\date{} 


% \pagestyle{myheadings}%         Set in Eng. 


\pagestyle{plain} %              Remove in Eng. 


\begin{document} 


% \thispagestyle{plain}%          Set in Eng. 


\maketitle 


\label{begin} 


 


Пусть $(\Omega,\frak A,P)$ --- вероятностное пространство, 


$(T,\frak B)$ --- 


измеримое пространство, $Y(t)$, $Y_n(t)$, $t\in T$, $n\in N$, --- 


независимые 


одинаково распределенные случайные процессы со значениями в 


банаховом пространстве $B$, определенные на $(\Omega,\frak A,P)$ и 


измеримые 


как функции двух аргументов относительно произведения 


$\sigma$-алгебр $\frak A\times\frak B$. Пусть $\xi$ и $\xi_n$, $n\in N$, 


--- независимые одинаково 


распределенные случайные величины со значениями в $(T,\frak B)$, 


определенные на другом вероятностном пространстве 


$(\Omega_1,\frak A_1,P_1)$. 


 


В [1]--[3] в случае, когда банахово пространство 


$B$ является действительной прямой $R$, изучается сходимость по 


распределению сумм и максимумов от случайных величин 


$Y_i(\xi_i(\omega_1))$ для почти всех $\omega_1\in\Omega_1$, сходимость 


для почти всех 


$\omega_1\in\Omega_1$ эмпирических процессов, случайных ломаных и 


случайных 


ступенчатых функций, определенных случайными величинами 


$Y_i(\xi_i(\omega_1))$. В [4] приведены обобщения результатов [1] на 


случайные процессы $Y_i$, принимающие значения в банаховом 


пространстве $B$. Здесь продолжаются эти исследования, 


изучается сходимость случайных ломаных, определенных нормированными 


суммами и нормированными билинейными формами от банаховозначных 


случайных элементов $Y_i(\xi_i(\omega_1))$ для почти всех 


$\omega_1\in\Omega_1$. 


 


Пусть $D$ --- подмножество вещественной прямой $R$, $L$ --- некоторое 


линейное пространство функций на $D$ со значениями в 


банаховом пространстве $B$, наделенное метризуемой топологией, 


$H(y)$, $H_n(y)$, $y\in D$, --- случайные процессы со значениями в 


$B$,  


почти все траектории которых принадлежат $L$. Говорят, что 


процессы $H_n(y)$ сходятся по распределению к процессу $H(y)$ в $L$ 


(и обозначают $H_n@>d>> H$, $n\to\infty$, в $L$) , если распределения 


случайных 


процессов $H_n$ сходятся слабо в $L$ к распределению случайного 


процесса $H$. Здесь в качестве пространства $L$ будем рассматривать 


пространство $C([0,1],B)$ непрерывных функций на 


отрезке $[0,1]$ со значениями в банаховом пространстве $B$. 


Символом $\deq$ будем обозначать равенство распределений 


случайных величин. 


 


Напомним ([5], с.\,54), что банахово пространство $B$ имеет тип~2, 


если существует такое $C>0$, что для всех $x_1,\dotsc,x_n\in B$, 


$n\in N$, 


$E\Big\|\sum\limits^n_{i=1}r_i x_i\Big\|^2\le 


C\Big(\sum\limits^n_{i=1}\|x_i\|^2\Big)$, где $r_i$ --- функции 


Радемахера (т.\,е.~случайные 


величины $r_i$ независимы и $P\{r_i=1\}=P\{r_i=-1\}=\frac{1}{2}$, 


$i\in N$). 


 


Пусть случайный процесс $Y$ и случайная величина $\xi$ такие, 


что $EY(t)=0$ для всех $t\in T$, ковариационная форма 


$V(x',y')=E_1Ex'(Y(\xi))y'(Y(\xi))$, $x',y'\in B'$, определяет 


гауссовский 


случайный элемент $G_{Y,\xi}$ в $B$ и найдется такой гауссовский 


$B$-значный случайный процесс $W_{Y,\xi}$, определенный на $[0,1]$ с 


независимыми приращениями, что для всех $x\in[0,1]$\; 


$W_{Y,\xi}(x)\deq x^{1/2}W_{Y,\xi}$. В этом случае будем говорить, что 


пара 


$(Y,\xi)$ определяет винеровский процесс $W_{Y,\xi}$. 


 


Рассмотрим случайные процессы в $B$, параметризованные 


элементами $\omega_1\in\Omega_1$, 


$$ 


S_n(x,\omega_1)=n^{-1/2}\bigg(\, 


\sum^{[nx]}_{i=1}Y_i(\xi_i(\omega_1))+\{nx\} 


Y_{[nx]+1}(\xi_{[nx]+1}(\omega_1))\bigg), 


$$ 


$0\le x\le 1$. Здесь через $[a]$ и $\{a\}$ обозначены целая и дробная 


части числа $a$. 


 


\begin{thm} Пусть банахово пространство $B$ имеет тип $2$, 


$EY(t)=0$ для всех $t\in T$, и $E_1E\|Y(\xi)\|^2<\infty$. Тогда пара 


$(Y,\xi)$ 


определяет винеровский процесс $W_{Y,\xi}$ и 


$S_n(x,\omega_1)@>d>> W_{Y,\xi}$, $n\to\infty$, 


в $C([0,1],B)$ для почти всех $\omega_1\in\Omega_1$. 


\end{thm} 


 


Пусть $H$ --- гильбертово пространство со скалярным произведением 


$(x,y)$, $x,y\in H$, случайные процессы $Y_i$ принимают 


значения в $H$. Рассмотрим случайные процессы, параметризованные 


элементами $\omega_1\in\Omega_1$, 


\begin{multline*} 


S^{\&}_n(x,\omega_1)=n^{-1}\bigg(\, 


\sum_{j<i\le[nx]}(Y_i(\xi_i(\omega_1)),Y_j(\xi_j(\omega_1)))+ \\ 


+\{nx\}\sum\begin{Sb} 


[nx]<i\le[nx]+1, \\ j<i\le[nx]+1 \end{Sb} 


(Y_i(\xi_i(\omega_1)),Y_j(\xi_j(\omega_1)))\bigg),\quad 


0\le x\le 1. 


\end{multline*} 


 


Обозначим через $I\in C[0,1]$ функцию $I(x)=x$, $x\in[0,1]$. 


 


\begin{thm} Пусть $B=H$ --- гильбертово пространство, 


$EY(t)=0$ для всех $t\in T$, и $E_1E\|Y(\xi)\|^2<\infty$. Тогда пара 


$(Y,\xi)$ 


определяет винеровский процесс $W_{Y,\xi}$ и 


$$ 


n^{-1}S^{\&}_n(x,\omega_1)@>d>> \tfrac{1}{2}((W_{Y,\xi}, 


W_{Y,\xi})-I),\ \; n\to\infty,\ \;\text{в }\; C[0,1] 


$$ 


для почти всех $\omega_1\in\Omega_1$. 


\end{thm} 


 


\begin{thebibliography}{9} 


 


\bibitem{1} 


Чупрунов А.Н. {\em О сходимости по распределению сумм и 


максимумов независимых случайных величин со случайными 


параметрами} // Liet. Matem. Rink. -- 1995. -- V.\,35. -- \No\,1. --


P.\,52--64. 


\bibitem{2} 


Чупрунов А.Н. {\em О сходимости по распределению эмпирических 


процессов, определенных независимыми случайными 


процессами} // Liet. Matem. Rink. -- 1995. -- V.\,35. -- \No\,2. 


-- P.\,171--180. 


\bibitem{3} 


Чупрунов А.Н. {\em О сходимости по распределению почти всех 


сумм независимых банаховозначных случайных элементов} // 


Изв. вузов. Математика. -- 1994. -- \No\,11. -- C.\,83--87. 


\bibitem{4} 


Чупрунов А.Н. {\em О сходимости случайных ломаных с 


нормировками типа Стьюдента} // Теория вероятн. и ее примен. 


-- 1996. -- T.\,51. -- \No\,4. -- C.\,914--919. 


\bibitem{5} 


Муштари Д.Х. {\em Вероятности и топологии в банаховых 


пространствах}. -- Казань: Изд-во Казанск. 


ун-та, 1989. -- 152 с. 


 


\end{thebibliography} 


 


\vskip 2cm 


 


\post{\it Казанский государственный}{\it Поступила} 


\post{\it университет}{13.11.1996} 


 


\label{end} 


\end{document} 


